\section{Introduction}
Latent Diffusion Models (LDMs) have become the dominant paradigm in image synthesis~\citep{rombach2022high, flux2024, qwenimage, sd3, seedream3, qwen-image-2.0}. These models typically employ a Variational Autoencoder (VAE) to project images into a compressed latent space for efficient diffusion modeling, where a widely adopted spatial compression ratio is 8 (denoted as $f8$). However, as the industry shifts toward native high-resolution synthesis, this standard ratio faces significant computational bottlenecks. Increasing the spatial compression ratio has thus become essential for computational efficiency, as the complexity of modern Diffusion Transformers (DiTs)~\citep{dit} scales quadratically with the number of latent tokens. Over the past few years, several advances have been achieved in this field, demonstrating the significant potential of high-compression VAEs~\citep{ltx, dcae, cosmos, dcae1.5}.


Despite these advances, a critical challenge exists: the inevitable trade-off between high compression ratio, reconstruction fidelity, and diffusability~\citep{diffusability}. Specifically, higher compression ratios often lead to severe degradation in reconstruction quality, particularly in text-rich scenarios where fine-grained detail is lost. While increasing the latent channel dimension can mitigate this information bottleneck, it frequently results in an over-complex and unstructured latent distribution, which significantly hinders the convergence and generative performance of downstream diffusion models~\citep{vavae, qiu2025image}.

In this work, we introduce Qwen-Image-VAE-2.0, a series of high-compression image VAEs ($f16$ \& $f32$), designed to overcome these challenges through improved architecture, comprehensive data engineering, and enhanced training strategy.

To address the challenge of reconstruction fidelity in high-compression regimes, we adopt an improved VAE architecture with Global Skip Connection (GSC), which establishes a global shortcut from pixels to latents, preserving fine-grained detail. Moreover, our design incorporates a higher latent dimensionality to alleviate the information bottleneck inherent in high-compression scenarios. On the data front, we scale our training corpus to billions of images and curate a specialized document collection (including academic papers, posters, slides, web pages, etc.) to enhance the reconstruction of text-rich images. Furthermore, we develop a synthetic pipeline that renders documents to provide dense supervisory signals for character-level reconstruction. Through these advancements, Qwen-Image-VAE-2.0 achieves state-of-the-art reconstruction performance, especially in text-rich scenarios, despite its high compression ratio.
To address the challenge of diffusability brought by high compression ratio and expanded latent dimension, we demonstrate that semantic alignment with intermediate features of DINOv2~\citep{dinov2} can effectively accelerate DiT convergence. Furthermore, we adopt a staged semantic alignment paradigm that transitions from strict semantic alignment to a balanced optimization of reconstruction and generation. Leveraging these techniques, Qwen-Image-VAE-2.0 achieves superior diffusability compared to existing high-compression VAEs, despite its large channel dimension.

To ensure computational efficiency, we leverage an asymmetric architecture that features a lightweight encoder to minimize encoding overhead during diffusion training. Additionally, we utilize an attention-free backbone to maintain high throughput even with ultra-high-resolution inputs. 

We conduct a comprehensive evaluation to validate the performance of Qwen-Image-VAE-2.0, focusing on two key aspects: reconstruction fidelity and latent space diffusability. For reconstruction fidelity, we evaluate it across general reconstruction tasks and introduce OmniDoc-TokenBench, a benchmark specifically targeting challenging scenarios like real-world text-rich document reconstruction. For latent space diffusability, we also perform extensive downstream DiT experiments to empirically verify it. The results demonstrate that Qwen-Image-VAE-2.0 not only achieves superior reconstruction fidelity, especially in text-rich scenarios despite high compression ratios, but also exhibits excellent latent space compatibility, facilitating rapid DiT convergence even with large latent dimension.

The key contributions of Qwen-Image-VAE-2.0 can be summarized as follows:
\begin{itemize}
\item \textbf{High-Compression VAE}: We introduce a suite of $f16$ and $f32$ image VAEs, providing a robust solution for efficient and native high-resolution image generation.
\item \textbf{Superior Reconstruction Performance:} Qwen-Image-VAE-2.0 achieves state-of-the-art reconstruction performance across multiple benchmarks. It maintains exceptional legibility in text-rich scenarios where traditional high-compression models typically fail.
\item \textbf{Enhanced Latent Diffusability:} Through a refined semantic alignment strategy, we demonstrate that large-channel VAEs can achieve excellent diffusability. This provides a promising solution to the tripartite trade-off between compression ratio, reconstruction fidelity, and diffusability.
\end{itemize}
